\documentclass[11pt]{amsart}
\usepackage{calc,amssymb,amsmath,ulem,stmaryrd,fullpage}
\usepackage{enumerate}
\usepackage{alltt}
\RequirePackage[dvipsnames,usenames]{color}
\let\mffont=\sf
\normalem
%\input{kmacros3.sty}
\input{xy}
\xyoption{all}
\usepackage{hyperref}
\usepackage{graphicx}
\usepackage[all,cmtip]{xy}
\usepackage{verbatim}
\usepackage[left=0.95in,top=0.95in,right=0.95in,bottom=0.95in]{geometry}
\newcommand{\tauCohomology}{T}
\newcommand{\FCohomology}{S}


\theoremstyle{theorem}
%\newtheorem*{mainthm*}{Main Theorem}

\newcommand{\note}[1]{\marginpar{\sffamily\tiny #1}}

\def\todo#1{\textcolor{Mahogany}%
{\sffamily\footnotesize\newline{\color{Mahogany}\fbox{\parbox{\textwidth-15pt}{#1}}}\newline}}

\renewcommand{\O}{\mathcal O}
\newcommand{\ie}{{\itshape i.e.} }
\newcommand{\roundup}[1]{\lceil #1 \rceil}
\newcommand{\rounddown}[1]{\lfloor #1 \rfloor}
\renewcommand{\to}[1][]{\xrightarrow{\ #1\ }}
\newcommand{\onto}[1][]{\protect{\xrightarrow{\ #1\ }\hspace{-0.8em}\rightarrow}}
\newcommand{\into}[1][]{\lhook \joinrel \xrightarrow{\ #1\ }}
%\newcommand{DBDual}[1]{{\underline{\omega}_{#1}}}
\newcommand{\DBDual}[1]{{{\uuline \omega}{}^{\mydot}_{#1}}}
\newcommand{\Tt}{{\mathfrak{T}}}
%\newcommand{\bn}{{\mathfrak{n}} }
\newcommand{\sol}[1]{ \vskip 6pt{\color{blue} {\bf Solution:} #1} \vskip 6pt}

\newcommand{\bR}{{\mathbb{R}}}
\newcommand{\va}{{\vec{a}}}
\newcommand{\vb}{{\vec{b}}}
\newcommand{\vzero}{{\vec{0}}}
\newcommand{\vi}{{\vec{i}}}
\newcommand{\vj}{{\vec{j}}}
\newcommand{\vk}{{\vec{k}}}
\newcommand{\vh}{{\vec{h}}}
\newcommand{\vr}{{\vec{r}}}
\newcommand{\vu}{{\vec{u}}}
\newcommand{\vv}{{\vec{v}}}
\newcommand{\vw}{{\vec{w}}}
\newcommand{\vx}{{\vec{x}}}
\newcommand{\vy}{{\vec{y}}}
\newcommand{\vz}{{\vec{z}}}
\newcommand{\vT}{{\vec{T}}}
\newcommand{\vN}{{\vec{N}}}
\newcommand{\vF}{{\vec{F}}}
\newcommand{\vG}{{\vec{G}}}
\newcommand{\bF}{\mathbb{F}}
\newcommand{\bQ}{\mathbb{Q}}
\newcommand{\bZ}{\mathbb{Z}}
\newcommand{\Inn}{\mathrm{Inn}}
\newcommand{\Aut}{\mathrm{Aut}}
\newcommand{\Ann}{\mathrm{Ann}}
\newcommand{\tor}{\mathrm{tor}}


\begin{document}
\title{Worksheet \#6 -- Math 6310\\Fall 2019}
\author{Due: Monday, November 18th}
\maketitle

\noindent
\vskip 9pt
\noindent
In this worksheet, we study classification of groups of small order.

Recall the following from worksheet \#2.  
\vskip 9pt
Let $N$ and $H$ be groups and let $\phi : H \to \text{Aut}(N)$ be a group homomorphism.  We get a left action of $H$ on $N$ by $h.n = (\phi(h))(n)$.
Set
\[
G = N \times H
\]
with the following binary operation:
\[
(n_1, h_1)(n_2, h_2) = (n_1 (h_1.n_2), h_1 h_2).
\]
This is called the \emph{semi-direct product of $N$ and $H$}, and is denoted by $N \rtimes H$.
\vskip 6pt
\noindent
{\bf Theorem.}  {\it Suppose $G$ is a group with $H$ an subgroup of $G$ and $N$ a normal subgroup of $G$.  Further suppose that $NH = G$ and that $N \cap H = \{ 1 \}$.  Let $\phi : H \to \mathrm{Aut}(N)$ be the map which conjugates by $h$.  In other words $\phi(h)(n) = hnh^{-1}$.  Then we have an isomorphism $N \rtimes H \cong G$.}
\vskip 3pt
We can use this to classify groups as follows.  Given some number $n$, suppose we can factor $n = ab$, with $\gcd(a,b) = 1$, where 
\begin{itemize}
    \item{} we can show (say by Sylow theorems) that groups $G$ of order $n$ have a normal subgroups $N$ of order $a$.
    \item{} and groups $G$ of order $n$ have other subgroups $H$ of order $b$.
\end{itemize}
Then any group of order $n$ must be isomorphic to a semidirect product $N \rtimes H$.  Thus to classify groups of order $n$, we need to classify such semidirect products.
\vskip 6pt
{\bf 1.}  Let $n = 10$.  Show that any group of order 10 has a normal subgroup $N$ of order 5 and a subgroup $H$ of order 2.  Classify all homomorphisms from $H \to \Aut(N)$.  How non-isomorphic groups of order 10 are there? (What happens when the homomorphism to $\Aut(N)$ is the trivial homomorphism?)
\newpage
\noindent
{\bf 2.}  Let $n = 20$.  Show that any group of order 20 has a normal subgroup $N$ of order 5 and a subgroup $H$ of order 4.  Classify groups of order 4 (they are all Abelian), and use that to classify groups of order 20.  
\vskip 3pt
\emph{Hint: }  There can be two different homomorphisms $\phi : H \to \Aut(N)$ that produce isomorphic semidirect products.  You'll need to figure out why that happens, at least for this particular example.  
\newpage
\noindent
{\bf 3.}  Classify groups of order 12.
\newpage
\noindent
{\bf 4.}  Show that if $n = p_1^{a_1} \dots p_l^{a_l}$ then $\Aut(\bZ/n\bZ)$ is isomorphic to 
\[
    U(\bZ/n\bZ) \cong U(\bZ/p_1^{a_1} \bZ) \oplus \dots \oplus U(\bZ/p_l^{a_l} \bZ)
\]
where $U(R)$ is the group of units of a ring $R$.  Additionally show that each $U(\bZ/p_i^{a_i} \bZ)$ is cyclic if $p_i \neq 2$.
\newpage
\noindent
{\bf 5.}  Classify groups of order 30 as follows.  Show that every group of order 30 has a subgroup of order 15 which is necessarily normal.  Next show/recall that groups of order 15 are cyclic.  Next show that $\Aut(\bZ/15) \cong \bZ/2 \times \bZ/4$ using the previous exercise.  Find the elements of order 2 in that group of automorphisms.  You can prove various groups you produce are nonisomorphic by looking at their centers.
\newpage
\noindent
{\bf 6.}  Classify all groups of order 18.
\newpage
\noindent
{\bf 7.}  Classify all groups of order 75.
\newpage
\noindent
{\bf 8.}  Classify all groups of order 70.

\end{document}

